\section{模型的评价}

\subsection{模型的优点}

\begin{itemize}
	\item \textbf{物理机理清晰}：基于非稳态热传导与水分扩散的偏微分方程建模，严格刻画了预热平衡与恒温干燥两个阶段，物性随含水率、温度变化的耦合关系处理得当，物理意义明确。
	\item \textbf{数值方法稳健高精度}：采用控制体积法空间离散与 Crank--Nicolson 隐式格式时间推进，无条件稳定、二阶精度，经网格收敛性检验验证了四位小数的输出精度，且天然满足质量守恒。
	\item \textbf{模型可扩展性强}：四个问题共用同一套求解框架，仅更换物性公式与边界条件即可从常物性推广到变物性、从固定区域推广到移动边界，体现了良好的统一性与复用性。
	\item \textbf{结果全面且可视化丰富}：给出了完整的温度场、水分场、干燥时间，并辅以剖面图、热力图、小提琴图、Sankey 流向图与 SHAP 贡献度图等多类型可视化，结论直观可信。
\end{itemize}

\subsection{模型的缺点}

\begin{itemize}
	\item \textbf{一维径向近似}：忽略端面传热传质，对长度有限、端面效应不可忽略的短粗药材会有一定偏差。
	\item \textbf{移动边界处理为近似}：问题四在固定网格上以等效传质系数近似移动边界，并忽略了收缩引起的对流输运项，属于工程近似而非严格动网格求解。
	\item \textbf{参数依赖经验公式}：$\rho$、$c_p$、$k$、$D$ 等均采用题目给定的经验公式，未考虑温度对热物性、含水率对收缩率等的更精细耦合。
\end{itemize}

\section{模型的改进与推广}

\subsection{模型的改进}

\begin{itemize}
	\item 采用二维轴对称 $(r,z)$ 模型，引入端面边界条件，以更精确地刻画短粗药材的三维效应。
	\item 对问题四采用 Landau 坐标变换或动网格方法严格处理收缩，显式加入材料点随收缩运动的对流项，提高移动边界模型的物理完备性。
	\item 引入蒸发潜热源项，建立热—湿—气三相耦合模型，并考虑干燥速率对表面边界条件的反馈，使模型更贴近真实干燥动力学。
\end{itemize}

\subsection{模型的推广}

本文的圆柱传热传质耦合模型及其数值求解框架，可推广到以下场景：

\begin{itemize}
	\item \textbf{其他农产品与中药材干燥}：茶叶、香菇、枸杞、人参等不同形状物料的干燥过程建模与工艺参数优化。
	\item \textbf{食品与化工颗粒干燥}：多孔介质、球形或片状物料的脱水过程，可沿用控制体积法统一求解。
	\item \textbf{逆向工艺设计}：以目标含水率与能耗为约束，结合本文灵敏度/SHAP 分析结果，反演最优烘房温度、湿度与通风风速，实现干燥工艺的数字化设计。
	\item \textbf{数值仿真平台}：将求解器封装为参数化仿真模块，用于构建“数字孪生”干燥车间，支持实时工艺监控与预测。
\end{itemize}
